\chapter{Introduzione}
\label{cap:introduzione}
\section{Descrizione dell'azienda}
Sync Lab s.r.L. è una azienda italiana attiva nell'abito \gls{ict}{}, specializzata nello sviluppo e consulenza IT dal 2002 con sedi a 
Milano, Roma, Napoli, Verona, Como, con più di 300 dipendenti. È una azienda orientata verso la \gls{Business Innovation}{}, finalizzata alla 
creazione di soluzioni innovative che abbracciano i nuovi paradigmi della trasformazione digitale.
Sync Lab possiede numerose certificazioni ISO LL-C per l'attestazione della 
qualità dei prodotti e servizi offerti. In particolare possiede le certificazioni 
ISO-9001 per la qualità dell'azienda, ISO-14001 per l'adozione quadro sistematico per l'integrazione delle pratiche a protezione dell'ambiente, ISO-27001 per la definizione di un \gls{sgsi}{}, ISO-45001 per l'adozione di un \gls{ohs}{}.
\\
Attualmente Sync Lab  lavora per più di 150 clienti diretti e finali, tra i più rilevanti ci sono nomi come: TIM, Trenitalia, PosteItaliane, UniCredit, ENI, ENEL, Vodafone, Fastweb.
\\
Inoltre è un'azienda che si pone come obiettivo quello di essere un punto di riferimento per i propri clienti nella realizzazione di prodotti e soluzioni innovative per diversi settori di mercato, come: Sanità, Industria, Energia, Telco, Finanza e Trasporti \& Logistica.\\
Lo spirito di Sync Lab è ampiamente rappresentato dal logo aziendale (Figura \ref{figure:logo_azienda}): un'onda che si propaga in modo circolare, che simboleggia la capacità di adattarsi e di evolversi in modo continuo.\\
\begin{figure}[htbp]  
\centering
    \includegraphics[width=0.5\textwidth]{images/introduzione/logo_azienda.png}
    \caption{Logo dell'aziedan Sync Lab s.r.L.}
    \label{figure:logo_azienda}
\end{figure}
\pagebreak
\section{Idea di fondo del progetto}
Oggigiorno la gestione e l'analisi di grandi moli di dati in tempo reale sta diventando fondamentale 
per le aziende che vogliono rimanere competitive sul mercato. \\ 
Per questo motivo è necessario utilizzare tecnologie e software che permettano di analizzare e archiviare 
i dati in tempo reale in modo efficiente e veloce. \\
D'altra parte è necessario anche che tali tecnologie siano in grado di scalare in modo verticale e orizzontale in base al carico 
di lavoro da sostenere, mantenendo sempre alte prestazioni e resilienza in caso di guasti. \\
Per questo motivo Sync Lab ha deciso d' investire in un progetto di ricerca e sviluppo che ha come obiettivo quello di creare 
una \gls{Data Pipeline}{} in grado di garantire le caratteristiche sopra descritte. \\
L'azienda ha già a disposizione un sistema di raccolta dati in real time, basato su Apache Kafka, che permette di ricevere dati da
diversi sistemi e applicazioni per poi inviarli ad un sistema di archiviazione. \\
Il progetto prevede l'inserimento di un sistema di \gls{Data Processing}{} basato su Apache Druid, che permetta di effettuare operazioni sui dati 
grezzi ricevuti da Apache Kafka, in modo da rendere più efficienti le successive operazioni di estrazione.\\
Particolarmente importante dovrà essere la fase di processing dei dati, in quanto dovrà permettere di eseguire operazioni di aggregazione
e trasformazione dei dati in modo efficiente e veloce riducendo al minimo i tempi di latenza e mantenendo alte prestazioni.\\
\subsection{Il ruolo dello stagista}
Lo stagista ha un ruolo fondamentale in tale tipologia di progetto, infatti è colui che porta uno spirito d'innovazione e consolida il valore aggiunto 
aziendale. \\
Le attività che costituiscono il percorso che lo stagista ha intrapreso sono state elencate all'interno di un \textit{Piano di lavoro}, concordato con il tutor aziendale, che ha lo scopo di guidare lo stagista durante il periodo di stage permettere al 
tutor aziendale di monitorare l'andamento delle attività delle attività svolte e di valutare il raggiungimento degli obiettivi.\\
Inoltre al termine del periodo di stage, sotto la supervisione del tutor aziendale è stata svolta una presentazione rivolta a tutti \textit{Stakeholder} aziendali, mirata a 
mostrare i risultati ottenuti con le tecnologie utilizzate e mettere in risalto le potenzialità del prototipo sviluppato. \\
Il progetto in sè fa parte di una rivoluzione tecnologica messa in atto da Sync Lab nel campo del \gls{Data Processing}{} e \gls{Data Analytics}{}.\\
\pagebreak
\section{Il progetto di stage}
\subsection{Descrizione del progetto}
Le attività descritte nel presente lavoro di stage illustrano la progettazione e lo sviluppo di un prototipo di una \gls{Data Pipeline}{} eseguibile
con Docker Compose, utilizzando Apache Druid come componente \gls{olap}{}.\\
Il prototipo finale riceve i dati da un sistema di raccolta basato su Apache Kafka eseguendo il \gls{Data Processing}{} con Apache Druid sui dati grezzi ricevuti.\\
L'obiettivo dello stage oltre a essere quello di sviluppare un prototipo funzionante, che soddisfi quanto richiesto, è anche quello di studiare e analizzare le funzionalità offerte dalle tecnologie utilizzate, in modo da poterne evidenziare i punti di forza e 
le differenze con le tecnologie tradizionali, utilizzate per il medesimo scopo. \\
\subsection{Obiettivi formativi}
In generale lo stage ha come obiettivo quello di far acquisire allo stagista concetti fondamentali riguardanti il contesto del prototipo sviluppato come: 
\begin{list}{*}{}
    \item Container technology;
    \item Apache Kafka e le Event Driven Architecture, design pattern publisher/subscriber;
    \item Column Based Database e la relazione/confronto con i classici Database relazionali;
    \item \gls{Middleware}{}, \gls{Data Pipeline}{}, le architetture distribuite, scalabili e resilienti.
\end{list}
\subsection{Risultati attesi e Obiettivi fissati}
I risultati attesi e gli obiettivi fissati per lo stage sono riportati nella Tabella \ref{tab:Tabella1}, con rispettivo identificativo, importanza e breve descrizione.
\\L'identificativo (riportato in breve con “ID”) è la sigla che identifica ogni requisito e rispetta la seguente notazione \textbf{[Importanza][Identificativo]}. \\
L’importanza è indicata dalla sigla \textbf{O} oppure \textbf{F} ad indicare rispettivamente un obiettivo
obbligatorio oppure facoltativo; mentre l’identificativo è un numero incrementale che
segnala in modo univoco l’obiettivo o il risultato in esame.\\

 \begin{table}[htbp]
    \centering
    \caption{Tabella degli obiettivi}    
    \label{tab:Tabella1}
    \begin{tabularx}{\textwidth}{|c|c|X|}
        \hline
        \textbf{ID} & \textbf{Importanza} & \textbf{Descrizione} \\\hline
        O1 & Obbligatorio & comprensione e definizione di una piccola \gls{Data Pipeline}{}  che  preveda il trattamento dei dati
        tramite Apache Kafka e Apache Druid \\\hline
        O2 & Obbligatorio & comprensione dei vantaggi e degli overhead  che le Event Driven Architecture portano con
        sé\\\hline
        O3 & Obbligatorio & comprensione del pattern publisher/subscriber \\\hline
        O4 & Obbligatorio & set-up di un cluster Apache Kafka in ambiente  containerizzato \\\hline
        O5 & Obbligatorio & gestione delle Time Series e dei Column-based Databases \\\hline
        O6 & Obbligatorio & comprensione delle differenze tra i database relazionali  classici e i Column-based Databases\\\hline
        O7 & Obbligatorio &comprensione dell’impiego e utilità dei \gls{Middleware}{} \\\hline
        F1 & Facoltativo & produzione di documentazione e un pacchetto di configurazione  dell’ambiente di sviluppo e
        esecuzione  della \gls{Data Pipeline}{}\\\hline
        F2 & Facoltativo & produzione di documentazione che riporti  le differenze  di performance  tra Apache Druid e altri
         database relazionali classici per alcune  operazioni \gls{olap}{} \\\hline
        F3 & Facoltativo & realizzazione di una presentazione che illustri l’architettura  sviluppata  a personale di settore o
        Stakeholder \\\hline
    \end{tabularx} 

\end{table}
\pagebreak
\pagebreak
\subsection{Analisi preventiva dei rischi}
Durante la fase di analisi iniziale del progetto di stage, sono stati individuati i seguenti rischi, cui si è cercato di porre rimedio con le azioni di mitigazione indicate. \\
\begin{enumerate}
    \item \textbf{Inesperienza tecnologica}: il progetto prevede l'utilizzo di tecnologie con cui lo stagista non ha mai avuto a che fare. \\
    \textbf{Rischio}: Medio.\\
    \textbf{Soluzione}: Per mitigare tale rischio, è stato previsto un periodo di ambientamento e formazione sulle tecnologie coinvolte, in modo da poter affrontare il progetto con maggiore consapevolezza.
    \item \textbf{Scelte errate nella progettazione dell'architettura}: il progetto prevede la progettazione di un'architettura complessa, con molte componenti, di natura differente che interagiscono tra di loro. \\
    \textbf{Rischio}: Alto.\\
    \textbf{Soluzione}: Per mitigare tale rischio, è stato previsto un periodo di analisi e progettazione dell'architettura, con il supporto del tutor aziendale, in modo da poter ovviare tale rischio.
    \item \textbf{Prestazioni insufficienti delle macchine a disposizione}: il progetto prevede l'impiego di tecnologie che richiedono un elevato dispendio di risorse. Tale fattore se non tenuto in considerazione potrebbe portare a risultati penalizzanti. \\
    \textbf{Rischio}: Alto.\\
    \textbf{Soluzione}: Per mitigare tale rischio, è stato previsto una configurazione di tali macchine in modo da poter sfruttare al meglio le risorse a disposizione.
\end{enumerate}    
\subsection{Obiettivi personali}
Nonostante la realizzazione del progetto sia l'obiettivo principale, il percorso di stage offre anche la possibilità di raggiungere una serie di obiettivi personali come: 


\begin{list}{*}{}
    \item imparare a utilizzare nuove tecnologie e strumenti legati ad architetture distribuite;
    \item comprendere i fattori da tenere in considerazione nella progettazione di un'architettura distribuita;
    \item comprendere i vantaggi e come suddividere il lavoro tra i componenti, in modo da poter lavorare in parallelo;
    \item imparare a lavorare in un team, condividendo le conoscenze e le esperienze;
    \item confrontarsi con persone del settore, per capire come si lavora in un'azienda.
\end{list}

\newpage
\pagestyle{empty}
\null % o \mbox{} o \phantom{X}

\newpage